% =====================================================================================
% WHITE PAPER one-page summary in a box (user 2026-11-12: "put the one-page idea in a box, in the first page; that page
% count is separate"). Numbers as in results.tex / discussion.tex. No functional / DFT details.
% =====================================================================================
\begin{tcolorbox}[colback=boxbg, colframe=boxframe, boxrule=0.8pt, arc=2pt, left=8pt, right=8pt, top=6pt, bottom=6pt,
  title={\textbf{At a glance}}, fonttitle=\large, coltitle=white]
\small
\textbf{The question.} Before making a perovskite film, can we predict how composition, annealing and substrate constraints will affect its structure, band gap and stability-related mechanisms? If so, we could identify promising processing conditions and prioritize experiments rather than relying primarily on trial and error.

\textbf{Why it matters.} The absorber is a polycrystalline film bonded to its substrate during the anneal; on
cooling the substrate leaves it under 50 to \SI{150}{MPa} of tension, reported to speed up degradation. In
\CsSnI{}, \MASnI{} and \FASnI{} the tilts that set the band gap change near device temperatures, and
black \CsSnI{} competes with a yellow phase.

\textbf{The short answer.} Processing strain barely moves the band gap or the overall rate of iodide migration, but it steers iodide: tension sends it along the film normal, towards the contacts. In \CsSnI{} it stabilises the photovoltaic black phase only slightly. Knowing what a process variable does not control also focuses optimisation.

\textbf{What we did.} With UrsaMLFF, our machine-learned force field, we took films of all three through an idealised anneal and cooling on a substrate, and computed band gap, iodide migration and black-phase stability.

\textbf{What we found.}
\begin{itemize}[leftmargin=12pt, itemsep=0pt, topsep=1pt]
  \item \textbf{The simulations reproduce the measured structures:} the tilt systems of all three compounds, tilt angles
    within \ang{0.3}, and the tilt transitions of the hybrids within the measured ranges.
  \item \textbf{The cation sets the tilt pattern} through the shape of its cage: FA (both ends hydrogen-bonding) prefers
    in-phase tilts by \SI{31}{meV} per formula unit, MA (one anchored end) antiphase tilts by 15 to \SI{17}{meV}, a rule for new cations.
  \item \textbf{Processing leaves 0.2 to 0.4\,\% tension} (50 to \SI{77}{MPa}, as measured in films), with every tilt pattern unchanged.
  \item \textbf{The band gap is not a processing variable:} the films open by 8 to \SI{19}{meV}, purely elastically,
    less than the 35 to \SI{67}{meV} by which the gap fluctuates at operating temperature.
  \item \textbf{Strain steers iodide instead of speeding it up:} per percent of tension, migration in \CsSnI{} becomes
    1.22 times faster along the normal and 0.87 times as fast in the plane, overall unchanged, as in MA films.
    Counter-intuitively, the stretch opens the paths towards the normal; the slight thinning narrows those in the plane.

  \item \textbf{Quantified, strain is no lever for \CsSnI{}:} residual tension gives only 2 to \SI{4}{meV} per
    formula unit of black-phase stabilisation (yellow is lower by 11 to \SI{19}{meV}) and raises migration towards the
    contacts by only about \SI{5}{\percent}.
\end{itemize}

\textbf{What it means for experiment and manufacturing.}
\begin{itemize}[leftmargin=12pt, itemsep=0pt, topsep=1pt]
  \item Substrate and anneal do not noticeably affect or deteriorate the band gap of tin iodide films; shifts of tens of meV point to
    composition or carrier concentration.
  \item Keep \CsSnI{} black by blocking yellow nucleation, not by strain; where iodide reaching the contacts limits
    lifetime, lower residual tension (cooler anneal, matched or hot-deposited transport layer) trims it slightly.
\end{itemize}

\textbf{How to test it.} Under tension, ion conduction in \CsSnI{} films should speed up across the film (about
\SI{20}{\percent} per percent) and slow along it (\SI{13}{\percent}) at unchanged activation energy; the gap of MA and FA films should not follow residual stress beyond about
\SI{15}{meV}; the yellow--black balance of \CsSnI{} should shift by only about \SI{0.2}{K} per MPa.

\textbf{Where this leads.} Today, the framework can model idealised annealing and substrate constraints for CsSnI\textsubscript{3}, MASnI\textsubscript{3} and FASnI\textsubscript{3}, predicting residual strain, structural response, band-gap changes and iodide migration in days. With grain boundaries, realistic interfaces and reactive atmospheres added and validated, the long-term goal is a virtual process line that helps researchers screen processing conditions, identify competing stability mechanisms and prioritize experiments before fabrication.
\end{tcolorbox}
